\documentclass[11pt,a4paper]{article}

\usepackage{amsmath,amsfonts,amssymb,amsthm}
\usepackage{mathtools}
\usepackage{graphicx}
\usepackage{hyperref}
\usepackage{geometry}
\geometry{margin=1in}

\title{\textbf{Multi-directional (J)ESPRIT with Random Base Points and PSF-Aware Weighting}}
\author{}
\date{}

\begin{document}

\maketitle

\section*{1. Problem setup}

We consider $r$ complex components with amplitudes $a_k \in \mathbb{C}$ and $d$-dimensional frequencies $\boldsymbol{\omega}_k \in [-\pi,\pi)^d$:
\begin{equation}
y[\mathbf{n}] = \sum_{k=1}^{r} a_k\, e^{j\langle \boldsymbol{\omega}_k, \mathbf{n}\rangle} + w[\mathbf{n}], 
\qquad \mathbf{n}\in\mathbb{Z}^d.
\end{equation}

If a known PSF with Fourier response $H(\cdot)$ is present, the Fourier-domain samples are multiplied by $H$, but \textbf{we never invert} $H$ (no pre-equalization).

\section*{2. Acquisition design (linear growth in $d$)}

\begin{itemize}
  \item Choose a small set of directions 
  $\mathcal{U} = \{\mathbf{u}_\ell\}_{\ell=1}^{M} \subset \mathbb{Z}^d$
  (e.g., canonical axes and a few short oblique directions).  
  Require $\operatorname{rank}(\mathbf{U}) = d$, where 
  $\mathbf{U} = [\mathbf{u}_1^\top; \dots; \mathbf{u}_M^\top] \in \mathbb{R}^{M\times d}$.
  \item For each direction $\mathbf{u}_\ell$, draw $S_\ell$ random base points 
  $\{\mathbf{p}_{s,\ell}\}_{s=1}^{S_\ell}$ with $S_\ell \gtrsim r$.
  \item For each $(\ell,s)$, sample the line
  \[
  \mathcal{L}_{\ell,s} = \{\,\mathbf{p}_{s,\ell} + n\mathbf{u}_\ell : n = 0,\dots,N_\ell-1\,\},
  \]
  with $N_\ell \ge r + 1$.
\end{itemize}

\section*{3. Per-direction model (multi-snapshot 1D)}

For direction $\ell$ and base point $s$:
\begin{equation}
z_{\ell,s}[n]
= H_\ell[n] \sum_{k=1}^{r} a_k\, e^{j\langle \boldsymbol{\omega}_k,\mathbf{p}_{s,\ell}\rangle}\, 
e^{j\langle \boldsymbol{\omega}_k,\mathbf{u}_\ell\rangle n} + w_{\ell,s}[n],
\end{equation}
where $H_\ell[n] = H(\mathbf{p}_{s,\ell}+n\mathbf{u}_\ell)$ (often independent of $s$).

Define:
\[
\varphi_{k,\ell} = \langle \boldsymbol{\omega}_k, \mathbf{u}_\ell\rangle, 
\qquad 
c_{k,s,\ell} = a_k\, e^{j\langle \boldsymbol{\omega}_k, \mathbf{p}_{s,\ell}\rangle}.
\]

Stacking the $S_\ell$ snapshots gives:
\begin{equation}
\mathbf{Z}_\ell = \underbrace{\mathbf{D}_\ell}_{\mathrm{diag}(H_\ell[0{:}N_\ell{-}1])}
\underbrace{\mathbf{V}_\ell}_{N_\ell \times r}
\underbrace{\mathbf{C}_\ell}_{r\times S_\ell} + \mathbf{W}_\ell,
\end{equation}
where 
$\mathbf{V}_\ell = [\mathbf{v}(\varphi_{1,\ell}),\dots,\mathbf{v}(\varphi_{r,\ell})]$ and
$\mathbf{v}(\phi) = [1, e^{j\phi}, \dots, e^{j\phi(N_\ell-1)}]^\top$.

\section*{4. Step 1 — Subspace estimation (multi-snapshot)}

For each $\ell$:
\begin{enumerate}
  \item Compute the covariance
  \(\widehat{\mathbf{R}}_\ell = \frac{1}{S_\ell}\mathbf{Z}_\ell \mathbf{Z}_\ell^H\).
  \item Extract the $r$ dominant eigenvectors 
  \(\mathbf{U}_{s,\ell} \in \mathbb{C}^{N_\ell\times r}\).
\end{enumerate}
No pre-equalization is used; the PSF weighting remains inside $\mathbf{D}_\ell$.

\section*{5. Step 2 — PSF-aware ESPRIT per direction}

Let $\mathbf{J}_\ell^{(1)}$ and $\mathbf{J}_\ell^{(2)}$ remove the last or first row, respectively.
Define
\[
\mathbf{\Lambda}_\ell^{(1)} = \mathbf{J}_\ell^{(1)}\mathbf{D}_\ell, \qquad
\mathbf{\Lambda}_\ell^{(2)} = \mathbf{J}_\ell^{(2)}\mathbf{D}_\ell.
\]

Estimate $\mathbf{\Phi}_\ell \in \mathbb{C}^{r\times r}$ using a PSF-weighted LS or TLS:
\begin{equation}
\widehat{\mathbf{\Phi}}_\ell
= \arg\min_{\mathbf{\Phi}}
\big\|\mathbf{W}_\ell(\mathbf{\Lambda}_\ell^{(2)}\mathbf{U}_{s,\ell}
- \mathbf{\Lambda}_\ell^{(1)}\mathbf{U}_{s,\ell}\mathbf{\Phi})\big\|_F,
\end{equation}
where $\mathbf{W}_\ell$ down-weights rows with small $|H_\ell[n]|$.
The eigenvalues of $\widehat{\mathbf{\Phi}}_\ell$ are
\[
\lambda_{k,\ell} = e^{j\widehat{\varphi}_{k,\ell}}, 
\qquad 
\widehat{\varphi}_{k,\ell} = \arg(\lambda_{k,\ell}).
\]

\section*{6. Step 3 — Joint ESPRIT (no pairing)}

All $\{\mathbf{\Phi}_\ell\}$ share the same eigenvectors.

\textbf{Option A (global subspace):}
build a single $\mathbf{U}_s$ from all directions and solve
\[
\mathbf{\Lambda}_\ell^{(2)}\mathbf{U}_s 
\approx \mathbf{\Lambda}_\ell^{(1)}\mathbf{U}_s\,\mathbf{\Phi}_\ell,
\quad \ell=1,\dots,M.
\]

\textbf{Option B (weak joint diagonalization):}
perform a joint diagonalization of the separate $\{\widehat{\mathbf{\Phi}}_\ell\}$:
\begin{equation}
\widehat{\mathbf{T}} = 
\arg\min_{\mathbf{T}}\sum_\ell \|\mathrm{offdiag}(\mathbf{T}^{-1}\widehat{\mathbf{\Phi}}_\ell \mathbf{T})\|_F^2.
\end{equation}

Each diagonal element corresponds to a consistent mode index $k$ across all $\ell$:
\[
\widehat{\varphi}_{k,\ell} = 
\arg\!\big((\widehat{\mathbf{T}}^{-1}\widehat{\mathbf{\Phi}}_\ell \widehat{\mathbf{T}})_{kk}\big).
\]
No pairing is needed.

\section*{7. Step 4 — Reconstruct $d$-D frequencies}

For each $k$, we have
$\widehat{\varphi}_{k,\ell} \approx \langle \boldsymbol{\omega}_k, \mathbf{u}_\ell\rangle \, (\mathrm{mod}\ 2\pi)$.

Let 
$\boldsymbol{\varphi}_k = [\widehat{\varphi}_{k,1},\dots,\widehat{\varphi}_{k,M}]^\top$.
Solve the small wrapped least-squares problem:
\begin{equation}
\min_{\boldsymbol{\omega}_k,\ \mathbf{m}\in\mathbb{Z}^{M}} 
\|\boldsymbol{\varphi}_k - \mathbf{U}\boldsymbol{\omega}_k - 2\pi \mathbf{m}\|_2^2.
\end{equation}

Practical two-step solution:
\begin{align}
\text{(a) unwrap:}& \quad \widetilde{\boldsymbol{\varphi}}_k \approx \text{unwrapped}(\boldsymbol{\varphi}_k),\\[2pt]
\text{(b) LS:}& \quad 
\widehat{\boldsymbol{\omega}}_k = 
(\mathbf{U}^\top \mathbf{U})^{-1}\mathbf{U}^\top \widetilde{\boldsymbol{\varphi}}_k.
\end{align}

\section*{8. Step 5 — Amplitude estimation (global weighted LS)}

Given $\{\widehat{\boldsymbol{\omega}}_k\}$, solve for $\mathbf{a} = [a_1,\dots,a_r]^\top$ by
\begin{equation}
\widehat{\mathbf{a}} 
= \arg\min_{\mathbf{a}}
\sum_{\mathbf{n}\in \cup_{\ell,s}\mathcal{L}_{\ell,s}} 
w(\mathbf{n})\,
\Big|z[\mathbf{n}] - 
H(\mathbf{n})\sum_{k=1}^{r} a_k e^{j\langle \widehat{\boldsymbol{\omega}}_k, \mathbf{n}\rangle}\Big|^2,
\end{equation}
with weights $w(\mathbf{n}) \propto |H(\mathbf{n})|^2$ or by discarding near-zero $H$.

\section*{9. Summary}

\begin{itemize}
  \item Sampling along random directions with multiple base points preserves 
  multi-snapshot structure.
  \item Each direction $\mathbf{u}_\ell$ is treated via a PSF-aware ESPRIT 
  (no $1/H$ inversion).
  \item A joint ESPRIT (global or weak JD) aligns the components automatically 
  without any pairing.
  \item $d$-D frequencies $\boldsymbol{\omega}_k$ are recovered by a small 
  wrapped least-squares, then amplitudes by global weighted LS.
\end{itemize}

\end{document}

